% IV. МЕТОДИКА И ЕКСПЕРИМЕНТАЛНИ РЕЗУЛТАТИ, част 2: ЧИСЛАТА.
% ВНИМАНИЕ: всяко число тук е измерено в средата, описана в 05_method.tex, и се
% възпроизвежда с посочената там команда. Клетка без измерване остава \TODO.
% NB: тук НЯМА \section. Файлът продължава раздела, открит в 05_method.tex.
\label{sec:results}

\subsection{Резултати}

\textbf{Вярност на извличането.}

Изпълнението е \code{ctest --test-dir build --output-on-failure}: 8 групи, 63 случая, всички
преминават, когато \code{clang++} е наличен; групата \code{extract} извежда \code{SKIP} и CTest я
отчита като пропусната, когато фронтендът липсва, а не като успешно преминала. Размерът на дървото и
на записания модел по-долу съдържат абсолютните пътища на
входните файлове, затова тези две колони се отместват с постоянна стойност при друг път; никоя
друга стойност не зависи от пътя.

\begin{table}[H]
\caption{Вярност на извличането по езикови конструкции}
\label{tab:extraction}
\centering
\setstretch{1.0}
\fontsize{10}{12}\selectfont
\begin{tabular}{@{}L{5.1cm}cL{5.1cm}c@{}}
\toprule
\textbf{Конструкция} & & \textbf{Конструкция} & \\
\midrule
Наследяване с достъп на базата & \cmark & Асоциативен контейнер & \cmark \\
Виртуално наследяване & \cmark & Статичен член, обхват на класа & \cmark \\
Множествено наследяване & \cmark & Чиста виртуална операция & \cmark \\
Член по стойност, композиция 1 & \cmark & Извеждане на интерфейс & \cmark \\
Притежаващ указател, композиция & \cmark & Шаблонен клас & \cmark \\
Споделен указател, агрегация & \cmark & Влагане в пространство от имена & \cmark \\
Указател без притежание & \cmark & Отсяване на подразбиращи се членове & \cmark \\
Член отправка, асоциация 1 & \cmark & Анотация за съхранение & \cmark \\
Контейнер от стойности, кратност * & \cmark & Отсяване по файл на произход & \cmark \\
Контейнер от указатели & \cmark & Породените заготовки са валиден C++ & \cmark \\
\bottomrule
\end{tabular}
\end{table}

Таблицата казва, че всяка изброена конструкция се извлича така, както е записано предварително. Тя
не казва, че покритието на езика е пълно: извън проверката остават приятелските декларации,
вложените шаблони на шаблони, наследяването от специализация на шаблон и всичко, което
препроцесорът е скрил при конкретния набор от ключове. Тези случаи не са проверени и затова не се
твърди нищо за тях.

\textbf{Цена на извличането.} Породените входове съдържат от 10 до 800 класа, всеки с по един представител на всяка форма на
член; съобщава се най-малкото от три повторения \tabref{tab:cost}.

\begin{table}[H]
\caption{Цена на извличането спрямо размера на единицата за превод}
\label{tab:cost}
\centering
\setstretch{1.0}
\fontsize{10}{12}\selectfont
\begin{tabular}{@{}rrrrrrrr@{}}
\toprule
\textbf{Класове} & \textbf{Изх. текст} & \textbf{Дърво} & \textbf{Възли} & \textbf{Елементи} &
\textbf{Фронтенд} & \textbf{Разбор} & \textbf{Обход} \\
 & \textbf{[B]} & \textbf{[B]} & & & \textbf{[s]} & \textbf{[s]} & \textbf{[s]} \\
\midrule
 10 &   3\,339 &      554\,127 &    3\,956 &    158 & 0.0631 & 0.0027 & 0.0004 \\
 25 &   8\,559 &      798\,297 &    5\,861 &    413 & 0.0688 & 0.0042 & 0.0010 \\
 50 &  17\,259 &   1\,206\,503 &    9\,036 &    838 & 0.0744 & 0.0067 & 0.0021 \\
100 &  34\,660 &   2\,024\,872 &   15\,386 & 1\,688 & 0.0932 & 0.0116 & 0.0055 \\
200 &  70\,360 &   3\,662\,672 &   28\,086 & 3\,388 & 0.1161 & 0.0215 & 0.0124 \\
400 & 141\,760 &   6\,945\,391 &   53\,486 & 6\,788 & 0.1790 & 0.0413 & 0.0323 \\
800 & 284\,560 &  13\,524\,380 &  104\,286 & 13\,588 & 0.2871 & 0.0815 & 0.0925 \\
\bottomrule
\end{tabular}
\end{table}

\begin{table}[H]
\caption{Двата действителни входа}
\label{tab:realinputs}
\centering
\setstretch{1.0}
\fontsize{10}{12}\selectfont
\begin{tabular}{@{}L{3.4cm}rrrrrr@{}}
\toprule
\textbf{Вход} & \textbf{Изх. текст} & \textbf{Дърво} & \textbf{Класи-} & \textbf{Връзки} &
\textbf{Модел} & \textbf{Общо} \\
 & \textbf{[B]} & \textbf{[B]} & \textbf{фикатори} & & \textbf{[B]} & \textbf{[s]} \\
\midrule
Анотиран пример & 3\,069 & 551\,549 & 6 & 8 & 6\,549 & 0.066 \\
Заглавни файлове на инструмента & 612 & 2\,952\,078 & 35 & 30 & 33\,366 & 0.108 \\
\bottomrule
\end{tabular}
\end{table}

От двете таблици следват четири неща. Първо, дървото расте линейно спрямо изходния текст, но с
коефициент около 48. Това е причината заместващите заглавни файлове да не са удобство, а условие: при
действителната стандартна библиотека същият коефициент дава стотици мегабайти. Второ, изпълнението на
фронтенда е най-скъпото до около 400 класа, но не расте пропорционално: при осемдесеткратно увеличение
на текста нараства 4.5 пъти, защото значителна част е постоянна цена за създаване на процес, докато
разборът на JSON расте практически линейно. Трето, обхождането до модел расте \term{по-бързо от
линейното}: между 400 и 800 класа входът се удвоява, а времето се увеличава 2.9 пъти, защото
разрешаването на име търси последователно в списък от всички класификатори. Поправката е очевидна,
речник вместо списък, и е записана като по-нататъшна работа, а не приложена, защото при действителните
входове разликата не е измерима. Четвърто, изходният текст на втория действителен вход е 612 байта,
защото файлът съдържа само включвания на деветте публични заглавни файла, а записаният модел е около 33
KB срещу близо 3 MB дърво. Това е количествената форма на твърдението, че моделът е абстракция: почти
цялото дърво е синтаксис, който диаграмата не показва.

\textbf{Запазване при пълен цикъл.}
\begin{table}[H]
\caption{Запазване на модела при пълен цикъл}
\label{tab:roundtrip}
\centering
\setstretch{1.0}
\fontsize{10}{12}\selectfont
\begin{tabular}{@{}L{3.7cm}L{2.2cm}rrrr@{}}
\toprule
\textbf{Вход} & \textbf{Път} & \textbf{Елементи} & \textbf{Запазени} & \textbf{Разлики} & \textbf{Дял} \\
\midrule
Анотиран пример & PlantUML & 41 & 41 & 0 & 1.000 \\
Анотиран пример & Mermaid  & 41 & 38 & 3 & 0.927 \\
Анотиран пример & схема    & 41 & 12 & 27 & 0.293 \\
Заглавни файлове & PlantUML & 226 & 226 & 0 & 1.000 \\
Заглавни файлове & Mermaid  & 226 & 225 & 1 & 0.996 \\
Заглавни файлове & схема    & 226 & 0 & 65 & 0.000 \\
\bottomrule
\end{tabular}
\end{table}

Пътят през PlantUML е без загуба и за двата входа. Това не беше вярно от самото начало: при първото
измерване се губеха имената на безименните параметри, защото излъчвателят ги записваше като
\code{: тип} и четецът прочиташе двоеточието обратно като част от типа. Дефектът беше открит именно от
този експеримент, а не от преглед на кода, и е поправен. Пътят през Mermaid губи по един елемент на
всяка операция, която е виртуална, но не е чиста: нотацията отбелязва операция като абстрактна или като
статична и няма трети знак. Загубата е свойство на нотацията, не на модела, и не подлежи на поправка.
Такива операции са три при анотирания пример и една при заглавните файлове, откъдето и разликата между
0.927 и 0.996, тоест делът зависи от това какъв дял от модела са виртуалните операции, а не от размера
му.

Пътят през схемата губи най-много, и това е очакваното: таблицата няма операции, няма видимост, няма
шаблони и няма класове, които не са обозначени за съхранение. При заглавните файлове се запазва нула,
защото нито един клас на инструмента не е обозначен. Нулата не е дефект, а вярно измерване на въпроса
колко от този модел изобщо може да бъде изразено като таблици. Общият извод е, че цикълът е затворен
през нотациите и е \term{частичен} през схемата: достатъчен, за да се отговори дали кодът и диаграмата
се разминават, и недостатъчен, за да се възстанови кодът от схемата. Второто и не се твърди.

\textbf{Откриваемост на внесено разминаване.}
\begin{table}[H]
\caption{Откриваемост на нарочно внесено разминаване}
\label{tab:divergence}
\centering
\setstretch{1.0}
\fontsize{10}{12}\selectfont
\begin{tabular}{@{}L{2.9cm}L{2.9cm}L{1.5cm}L{2.1cm}L{2.4cm}@{}}
\toprule
\textbf{Вид на промяната} & \textbf{Очакван вид съобщение} & \textbf{Открита} & \textbf{Общо съобщения} & \textbf{Извън очакваното} \\
\midrule
Премахнат атрибут & добавен член      & \cmark & 1 & 0 \\
Добавена операция & премахнат член    & \cmark & 1 & 0 \\
Сменена видимост  & променен член     & \cmark & 1 & 0 \\
Премахнат базов клас & добавена връзка & \cmark & 1 & 0 \\
Сменена кратност  & променена връзка  & \cmark & 2 & 1 \\
\bottomrule
\end{tabular}
\end{table}

И петте внесени разминавания са открити, при нула съобщения, които не отговарят на действително
разминаване. Второто число е важното: инструмент, който съобщава за всичко, би постигнал същата
откриваемост и би бил безполезен. Четири от петте случая дават точно по едно съобщение. Петият дава
две: при смяна на кратността от \code{0..1} на \code{1} се съобщава и за променена връзка, и за
променен атрибут. Второто не е грешно, защото кратността се пази на две места, върху връзката и върху
атрибута, който я е породил, и промяната на диаграмата променя и двете, тоест едно редактиране води
до две верни съобщения. Наблюдението обаче сочи истинско проектно решение: кратността е
\term{дублирана} в модела. Дублирането е удобно за излъчвателите и неудобно за проверяващия модул,
който брои една промяна два пъти; отстраняването му е въпрос на избор кой носител е първичен, и е
записано като по-нататъшна работа.
